\usepackage{tcolorbox}
\usepackage{minted}
% JetBrains Mono as the monospace/code font, embedded as TTF via \pdfmapline
% (same mechanism as NVIDIA Sans in the class). No package dependency, works
% under pdflatex on Overleaf and arXiv.
\ProvidesFile{t1JetBrainsMono.fd}
\DeclareFontFamily{T1}{JetBrainsMono}{}
\DeclareFontShape{T1}{JetBrainsMono}{m}{n}{<-> nvidia_release_template/JetBrainsMono-Font-TTF/JetBrainsMono_Rg}{}
\DeclareFontShape{T1}{JetBrainsMono}{b}{n}{<-> nvidia_release_template/JetBrainsMono-Font-TTF/JetBrainsMono_Bd}{}
\DeclareFontShape{T1}{JetBrainsMono}{m}{it}{<-> nvidia_release_template/JetBrainsMono-Font-TTF/JetBrainsMono_It}{}
\DeclareFontShape{T1}{JetBrainsMono}{b}{it}{<-> nvidia_release_template/JetBrainsMono-Font-TTF/JetBrainsMono_BdIt}{}
\DeclareFontShape{T1}{JetBrainsMono}{bx}{n}{<->ssub * JetBrainsMono/b/n}{}
\DeclareFontShape{T1}{JetBrainsMono}{bx}{it}{<->ssub * JetBrainsMono/b/it}{}
\DeclareFontShape{T1}{JetBrainsMono}{m}{sl}{<->ssub * JetBrainsMono/m/it}{}
\DeclareFontShape{T1}{JetBrainsMono}{b}{sl}{<->ssub * JetBrainsMono/b/it}{}
\DeclareFontShape{T1}{JetBrainsMono}{m}{sc}{<->ssub * JetBrainsMono/m/n}{}

\pdfmapline{+JetBrainsMono_Rg < nvidia_release_template/JetBrainsMono-Font-TTF/JetBrainsMono-Regular.ttf <T1-WGL4.enc}
\pdfmapline{+JetBrainsMono_Bd < nvidia_release_template/JetBrainsMono-Font-TTF/JetBrainsMono-Bold.ttf <T1-WGL4.enc}
\pdfmapline{+JetBrainsMono_It < nvidia_release_template/JetBrainsMono-Font-TTF/JetBrainsMono-Italic.ttf <T1-WGL4.enc}
\pdfmapline{+JetBrainsMono_BdIt < nvidia_release_template/JetBrainsMono-Font-TTF/JetBrainsMono-BoldItalic.ttf <T1-WGL4.enc}
\renewcommand{\ttdefault}{JetBrainsMono}
\tcbuselibrary{minted,skins,breakable}
\usepackage{mathrsfs}
\usepackage{amsmath}
\usepackage{physics}
\usepackage{mathtools}
\usepackage{dsfont} % for mathds
\usepackage{float}



\newcommand{\algcomshort}[1]{\textcolor{blue}{{\hfill $\triangleright$ { #1}}}}
\usepackage{subcaption}
\usepackage{graphicx}
\usepackage{geometry} % For adjusting margins if needed


% for tables.
\usepackage{booktabs} % For \toprule, \midrule, \bottomrule
\usepackage{multirow} % For \multirow command
\usepackage{array}    % For better column formatting
\usepackage{tcolorbox}

% for makecell
\usepackage{makecell} % For multi-line headers

\definecolor{codeBg}{HTML}{F7F8FA}
\definecolor{codeFrame}{HTML}{D9DEE7}
\definecolor{codeTitleBg}{HTML}{2F3747}
\definecolor{codeTitleText}{HTML}{FFFFFF}
\definecolor{codeLine}{HTML}{7A8494}
% --- VS Code Light syntax palette ---
\definecolor{codeKeyword}{HTML}{0000FF}
\definecolor{codeString}{HTML}{A31515}
\definecolor{codeComment}{HTML}{008000}
\definecolor{codeNumber}{HTML}{098658}
\definecolor{codeType}{HTML}{267F99}
\definecolor{codeFunc}{HTML}{795E26}
\makeatletter
\newcommand{\appendixcodefontsize}{%
  \footnotesize
  \fontsize{0.8\dimexpr\f@size pt\relax}{0.96\dimexpr\f@size pt\relax}\selectfont
}
\makeatother
\renewcommand{\theFancyVerbLine}{\textcolor{codeLine}{\scriptsize\arabic{FancyVerbLine}}}
\makeatletter
% VS Code Light palette: recolor Pygments tokens and drop default italics.
\newcommand{\ApplyCodePalette}{%
  % comments (green, upright)
  \@namedef{PYG@tok@c}{\def\PYG@tc####1{\textcolor{codeComment}{####1}}}%
  \@namedef{PYG@tok@ch}{\def\PYG@tc####1{\textcolor{codeComment}{####1}}}%
  \@namedef{PYG@tok@cm}{\def\PYG@tc####1{\textcolor{codeComment}{####1}}}%
  \@namedef{PYG@tok@cp}{\def\PYG@tc####1{\textcolor{codeComment}{####1}}}%
  \@namedef{PYG@tok@cpf}{\def\PYG@tc####1{\textcolor{codeComment}{####1}}}%
  \@namedef{PYG@tok@c1}{\def\PYG@tc####1{\textcolor{codeComment}{####1}}}%
  \@namedef{PYG@tok@cs}{\def\PYG@tc####1{\textcolor{codeComment}{####1}}}%
  % keywords (blue, bold preserved from the default style)
  \@namedef{PYG@tok@k}{\let\PYG@bf\textbf\def\PYG@tc####1{\textcolor{codeKeyword}{####1}}}%
  \@namedef{PYG@tok@kc}{\let\PYG@bf\textbf\def\PYG@tc####1{\textcolor{codeKeyword}{####1}}}%
  \@namedef{PYG@tok@kd}{\let\PYG@bf\textbf\def\PYG@tc####1{\textcolor{codeKeyword}{####1}}}%
  \@namedef{PYG@tok@kn}{\let\PYG@bf\textbf\def\PYG@tc####1{\textcolor{codeKeyword}{####1}}}%
  \@namedef{PYG@tok@kp}{\let\PYG@bf\textbf\def\PYG@tc####1{\textcolor{codeKeyword}{####1}}}%
  \@namedef{PYG@tok@kr}{\let\PYG@bf\textbf\def\PYG@tc####1{\textcolor{codeKeyword}{####1}}}%
  % types and builtins (teal); class bold preserved
  \@namedef{PYG@tok@kt}{\def\PYG@tc####1{\textcolor{codeType}{####1}}}%
  \@namedef{PYG@tok@nc}{\let\PYG@bf\textbf\def\PYG@tc####1{\textcolor{codeType}{####1}}}%
  \@namedef{PYG@tok@nb}{\def\PYG@tc####1{\textcolor{codeType}{####1}}}%
  % functions and decorators (ochre)
  \@namedef{PYG@tok@nf}{\def\PYG@tc####1{\textcolor{codeFunc}{####1}}}%
  \@namedef{PYG@tok@fm}{\def\PYG@tc####1{\textcolor{codeFunc}{####1}}}%
  \@namedef{PYG@tok@nd}{\def\PYG@tc####1{\textcolor{codeFunc}{####1}}}%
  % strings (dark red)
  \@namedef{PYG@tok@s}{\def\PYG@tc####1{\textcolor{codeString}{####1}}}%
  \@namedef{PYG@tok@sd}{\def\PYG@tc####1{\textcolor{codeString}{####1}}}%
  \@namedef{PYG@tok@s1}{\def\PYG@tc####1{\textcolor{codeString}{####1}}}%
  \@namedef{PYG@tok@s2}{\def\PYG@tc####1{\textcolor{codeString}{####1}}}%
  \@namedef{PYG@tok@sb}{\def\PYG@tc####1{\textcolor{codeString}{####1}}}%
  \@namedef{PYG@tok@se}{\def\PYG@tc####1{\textcolor{codeString}{####1}}}%
  \@namedef{PYG@tok@si}{\def\PYG@tc####1{\textcolor{codeString}{####1}}}%
  \@namedef{PYG@tok@sr}{\def\PYG@tc####1{\textcolor{codeString}{####1}}}%
  \@namedef{PYG@tok@ss}{\def\PYG@tc####1{\textcolor{codeString}{####1}}}%
  \@namedef{PYG@tok@dl}{\def\PYG@tc####1{\textcolor{codeString}{####1}}}%
  % numbers (green)
  \@namedef{PYG@tok@m}{\def\PYG@tc####1{\textcolor{codeNumber}{####1}}}%
  \@namedef{PYG@tok@mi}{\def\PYG@tc####1{\textcolor{codeNumber}{####1}}}%
  \@namedef{PYG@tok@mf}{\def\PYG@tc####1{\textcolor{codeNumber}{####1}}}%
  \@namedef{PYG@tok@mh}{\def\PYG@tc####1{\textcolor{codeNumber}{####1}}}%
  \@namedef{PYG@tok@mo}{\def\PYG@tc####1{\textcolor{codeNumber}{####1}}}%
  \@namedef{PYG@tok@mb}{\def\PYG@tc####1{\textcolor{codeNumber}{####1}}}%
  \@namedef{PYG@tok@il}{\def\PYG@tc####1{\textcolor{codeNumber}{####1}}}%
}
\makeatother

\DeclareCaptionType{codesnippet}[Code Snippet][List of Code Snippets]
\captionsetup[codesnippet]{
  labelfont={},
  labelsep=colon,
  justification=raggedright,
  singlelinecheck=false,
  skip=3pt,
  hypcap=false
}
\crefname{codesnippet}{code snippet}{code snippets}
\Crefname{codesnippet}{Code Snippet}{Code Snippets}

\newtcblisting{appendixcode}[2][]{%
  enhanced,
  breakable,
  width=\linewidth,
  listing only,
  listing engine=minted,
  minted language=#2,
  minted options={
    fontsize=\appendixcodefontsize,
    linenos,
    numbersep=6pt,
    xleftmargin=1.8em,
    breaklines,
    autogobble,
    tabsize=4,
    formatcom=\ApplyCodePalette
  },
  colback=codeBg,
  colframe=codeFrame,
  boxrule=0.45pt,
  arc=0mm,
  outer arc=0mm,
  left=0mm,
  right=1mm,
  top=1mm,
  bottom=1mm,
  boxsep=1mm,
  before skip=4pt,
  #1
}

%%%%%%%%%%%%%%%%%%%%%%%%%%%%%%%%%%%%%%%%%%%%%%%%%%%%%%%%%%%%%%%%%
%%%%%%%%%%%%%%%%%%%%  font color setting  %%%%%%%%%%%%%%%%%%%%%%%
%%%%%%%%%%%%%%%%%%%%%%%%%%%%%%%%%%%%%%%%%%%%%%%%%%%%%%%%%%%%%%%%%
\newcommand{\red}[1]{\textcolor{red}{#1}}
\newcommand{\blue}[1]{\textcolor{blue}{#1}}
\newcommand{\orange}[1]{\textcolor{orange}{#1}}
\newcommand{\green}[1]{\textcolor{ForestGreen}{#1}}
\newcommand{\purple}[1]{\textcolor{purple}{#1}}
\newcommand{\pink}[1]{\textcolor{pink}{#1}}
\newcommand{\ysa}[1]{\textcolor{green}{#1}}
%%%%%%%%%%%%%%%%%%%%%%%%%%%%%%%%%%%%%%%%%%%%%%%%%%%%%%%%%%%%%%%%%
% Project website. Default = public URL. Anonymous builds (main.tex,
% arxiv.tex) override \websiteurl after inputting this file.
\newcommand{\websiteurl}{enpire-research.github.io}
\newcommand{\website}{\href{https://\websiteurl}{\websiteurl}}
%%%%%%%%%%%%%%%%%%%%%%%%%%%%%%%%%%%%%%%%%%%%%%%%%%%%%%%%%%%%%%%%%
%\usepackage{hyperref}
%\hypersetup{colorlinks,linkcolor={ForestGreen},citecolor={blue},urlcolor={blue}}
%%%%%%%%%%%%%%%%%%%%%%%%%%%%%%%%%%%%%%%%%%%%%%%%%%%%%%%%%%%%%%%%%
%%%%%%%%%%%%%%%%%%%%%%%%%%%%%%%%%%%%%%%%%%%%%%%%%%%%%%%%%%%%%%%%%

%%%%%%%%%%%%%%%%%%%%%%%%%%%%%%%%%%%%%%%%%%%%%%%%%%%%%%%%%%%%%%%%%
%%%%%%%%%%%%%%%%%%%%%% mathematical commands %%%%%%%%%%%%%%%%%%%%
%%%%%%%%%%%%%%%%%%%%%%%%%%%%%%%%%%%%%%%%%%%%%%%%%%%%%%%%%%%%%%%%%
\DeclareMathAlphabet{\mathcal}{OMS}{cmsy}{m}{n}
\SetMathAlphabet{\mathcal}{bold}{OMS}{cmsy}{b}{n}
\newcommand{\bigO}{\mathcal{O}}
\newcommand{\bigOmega}{\mathcal{O}}
\newcommand{\cT}{\mathcal{T}}
\newcommand{\cX}{\mathcal{X}}

\newcommand{\T}{\mathbb{T}}
\newcommand{\B}{\mathbf{B}}
\renewcommand{\O}{\mathbb{O}}
\newcommand{\Z}{\mathbb{Z}}
\newcommand{\R}{\mathbb{R}}
\newcommand{\lt}{\ensuremath <}
\newcommand{\gt}{\ensuremath >}

\newcommand{\D}{\mathscr{D}} 
\renewcommand{\r}{\boldsymbol{\rho}}
\newcommand{\p}{\mathbf{p}}
\newcommand{\BS}{\Delta(\mathscr{S})}  %Belief State
\newcommand{\E}{\mathbb{E}}
\renewcommand{\|}{\Bigg|}
\renewcommand{\[}{\left[}
\renewcommand{\]}{\right]}
\renewcommand{\(}{\left(}
\renewcommand{\)}{\right)}
\newcommand{\x}{\mathbf{x}}
\newcommand{\s}{\mathbf{s}}
\renewcommand{\a}{\mathbf{a}}
\newcommand{\m}{\mathbf{m}}
\newcommand{\n}{\mathbf{n}}
\newcommand{\e}{\mathbf{e}}
\renewcommand{\v}{\mathbf{v}}
\renewcommand{\u}{\mathbf{u}}
\renewcommand{\P}{\mathbb{P}}
\newcommand{\y}{\mathbf{y}}
\renewcommand{\H}{\mathbf{H}}
\newcommand{\Regret}{\text{Regret}}
\renewcommand{\wp}{\text{with probability at least $1-\delta$}}
\newcommand{\abbr}{\overset{abbr}{=}}
%spaces
\newcommand{\sS}{\mathscr{S}}
\newcommand{\sA}{\mathscr{A}} 
\newcommand{\sO}{\mathscr{O}} 
\newcommand{\sU}{\mathscr{U}}
\newcommand{\sF}[1]{\mathscr{F}_{#1}}
\newcommand{\sT}[1]{\mathcal{T}_{#1}} 
%random variables
\renewcommand{\S}{\mathbf{S}}
\newcommand{\A}{\mathbf{A}}
\renewcommand{\O}{\mathbf{O}} 
\newcommand{\F}{{\boldsymbol{F}}}
\newcommand{\f}{{f}}
\newcommand{\btau}{\boldsymbol{\tau}}

%matrices
\newcommand{\bT}{\mathbb{T}}
\newcommand{\bO}{\mathbb{O}}
\newcommand{\bM}{\mathbb{M}}
\newcommand{\bL}{\Lambda}
\newcommand{\K}{\mathcal{K}}
\renewcommand{\b}{{\vec{b}}}
%Shortcuts for operators
% \DeclareMathOperator*{\argmax}{argmax}
% \DeclareMathOperator*{\argmin}{argmin}
\newcommand{\argmax}[1]{\mathrm{argmax}}
\newcommand{\argmin}[1]{\mathrm{argmin}}

\newcommand{\Min}[1]{\underset{ #1 }{\min \ }}
\newcommand{\sgn}{\text{sgn}}
\newcommand{\Pm}[1]{\mathbb{P}_{\mathcal{P}}\left\{ #1 \right\}}
%probability given the model
\renewcommand{\Pr}[1]{\mathbb{P}_{\mathcal{P}}^\pi\left\{ #1 \right\}} %probability given the model and the policy
\newcommand{\diag}[1]{\text{diag}\left(#1\right)}%diagonal matrix
\renewcommand{\over}[1]{\overset{#1}{=}}
\newcommand{\muu}[1]{\vec{\mu}_{#1}}
\newcommand{\V}{\mathsf{V}}
%sign function
\newcommand{\sign}[1]{\sgn{#1}}
%math symbols
% \newcommand{\abs}[1]{\left\lvert#1\right\rvert}
\newcommand{\algcom}[1]{\textcolor{blue}{ #1}}
% Define the \norm command




%%%%%%%%%%%%%%%%%%%%%%%%%%%%%%%%%%%%%%%%%%%%%%%%%%%%%%%%%%%%%%%%%
%%%%%%%%%%%%%%%%%%%%  theorem style  %%%%%%%%%%%%%%%%%%%%%%%
%%%%%%%%%%%%%%%%%%%%%%%%%%%%%%%%%%%%%%%%%%%%%%%%%%%%%%%%%%%%%%%%%
\usepackage{amsthm}
\theoremstyle{plain}
\newtheorem{theorem}{Theorem}[section]
\newtheorem{proposition}[theorem]{Proposition}
\newtheorem{lemma}[theorem]{Lemma}
\newtheorem{corollary}[theorem]{Corollary}
\theoremstyle{definition}
\newtheorem{definition}[theorem]{Definition}
\newtheorem{assumption}[theorem]{Assumption}
\theoremstyle{remark}
\newtheorem{remark}[theorem]{Remark}
\usepackage{cancel}
%%%%%%%%%%%%%%%%%%%%%%%%%%%%%%%%%%%%%%%%%%%%%%%%%%%%%%%%%%%%%%%%%
%%%%%%%%%%%%%%%%%%%% algorithm   %%%%%%%%%%%%%%%%%%%%%%%
%%%%%%%%%%%%%%%%%%%%%%%%%%%%%%%%%%%%%%%%%%%%%%%%%%%%%%%%%%%%%%%%%
\usepackage{algorithm}
% \usepackage{algorithmic}
\usepackage{algpseudocode}
